% =====================================================================
% Rozdział 6: Zagadnienia bezpieczeństwa i prywatności
% =====================================================================

\chapter{Zagadnienia bezpieczeństwa i prywatności}
\label{chap:bezpieczenstwo}

\section{Zakres odpowiedzialności modelu}
Model LEM definiuje wyłącznie semantykę informacji o lokalizacji oraz zasady jej interpretacji. LEM NIE DEFINIUJE mechanizmów uwierzytelniania, autoryzacji, kontroli dostępu, szyfrowania, ochrony transmisji, retencji danych ani zasad zarządzania informacjami lokalizacyjnymi.

Implementacje wykorzystujące LEM POWINNY zapewniać mechanizmy ochrony danych odpowiednie do wymagań środowiska, w którym model jest stosowany. Model LEM nie definiuje ani nie weryfikuje spełnienia tego wymagania — pozostaje ono w gestii implementacji i środowiska, podobnie jak inne wymagania pozamodelowe określone w sekcji 2.5.5.

Bezpieczeństwo informacji lokalizacyjnej jest zależne od sposobu jej pozyskania, przechowywania, przetwarzania oraz udostępniania przez systemy wykorzystujące model LEM.

\section{Ochrona integralności informacji}
Informacja lokalizacyjna reprezentowana zgodnie z LEM MUSI zachowywać zgodność z interpretacją kanoniczną modelu.

Implementacja NIE MOŻE modyfikować znaczenia komponentów lokalizacji w sposób prowadzący do zmiany ich interpretacji określonej w Rozdziale 2. Każda modyfikacja informacji lokalizacyjnej, która narusza zasady semantyczne określone w modelu LEM, unieważnia zgodność tych danych z modelem.

Mechanizmy zapewniające integralność danych, w tym kontrola zmian, podpisywanie, sumy kontrolne lub inne metody ochrony informacji, pozostają odpowiedzialnością implementacji.

\section{Separacja lokalizacji od identyfikacji}
Model LEM opisuje informację o fizycznej lokalizacji niezależnie od identyfikacji obiektu, użytkownika lub systemu.

LEM NIE DEFINIUJE:
\begin{itemize}
    \item mechanizmów identyfikacji obiektów,
    \item mechanizmów identyfikacji użytkowników,
    \item sposobów powiązania lokalizacji z konkretną osobą lub podmiotem,
    \item zasad zarządzania tożsamością.
\end{itemize}

Powiązanie informacji lokalizacyjnej z identyfikatorami obiektów lub osób jest wykonywane poza modelem LEM i podlega zasadom obowiązującym w danym systemie. Identyfikatory nadane przez systemy zewnętrzne mogą być wykorzystywane przez implementacje LEM wyłącznie jako sposób odwołania się do elementu opisu lokalizacji (zob. A.9) — takie wykorzystanie nie stanowi mechanizmu identyfikacji w rozumieniu modelu LEM ani nie osłabia zasady separacji określonej w niniejszej sekcji.

Implementacje MUSZĄ uwzględniać, że połączenie informacji lokalizacyjnej z innymi danymi może powodować powstanie dodatkowych wymagań bezpieczeństwa lub prywatności.

\section{Zaufanie do źródła informacji lokalizacyjnej}
Model LEM nie określa metod pozyskiwania informacji lokalizacyjnej ani sposobów oceny wiarygodności źródła danych.

LEM NIE DEFINIUJE:
\begin{itemize}
    \item technologii pomiaru lokalizacji,
    \item algorytmów określania położenia,
    \item metod oceny dokładności lub jakości pomiaru,
    \item mechanizmów potwierdzania pochodzenia danych.
\end{itemize}

Implementacja wykorzystująca LEM MOŻE przechowywać dodatkowe informacje dotyczące pochodzenia lub jakości danych, jednak informacje te nie stanowią części podstawowej semantyki modelu LEM, chyba że zostaną określone przez odpowiedny profil zgodny z zasadami Rozdziału 4.

\section{Profile bezpieczeństwa}
Wymagania bezpieczeństwa specyficzne dla określonych środowisk MOGĄ być definiowane poprzez profile bezpieczeństwa lub profile zastosowań zgodne z zasadami określonymi w Rozdziale 4.

Profil bezpieczeństwa:
\begin{itemize}
    \item NIE MOŻE zmieniać znaczenia komponentów modelu LEM,
    \item NIE MOŻE definiować alternatywnej interpretacji informacji lokalizacyjnej,
    \item MOŻE określać wymagania dotyczące ochrony, dostępu lub sposobu wykorzystania informacji.
\end{itemize}

Profile bezpieczeństwa pozostają ograniczeniami dotyczącymi implementacji lub zastosowania modelu, a nie rozszerzeniem semantyki LEM.

\section{Prywatność informacji lokalizacyjnej}
Model LEM jest neutralny względem klasyfikacji prawnej lub organizacyjnej informacji lokalizacyjnej.

LEM NIE OKREŚLA, czy określona informacja lokalizacyjna stanowi dane osobowe, informację poufną lub inną kategorię danych chronionych.

Ocena charakteru informacji oraz zastosowanie odpowiednich mechanizmów ochrony pozostają odpowiedzialnością systemów wykorzystujących model LEM.

Implementacje MUSZĄ uwzględniać wymagania dotyczące prywatności wynikające z obowiązujących regulacji, polityk organizacyjnych oraz charakteru przetwarzanych danych. W szczególności implementacje POWINNY uwzględniać, że nawet ograniczony podzbiór komponentów LEM (np. Reference Space i Relative Position bez pełnego kontekstu hierarchicznego) może, w zależności od kontekstu, stanowić informację o istotnym znaczeniu dla prywatności (zob. 2.4.4).

\section{Minimalizacja ujawnianej informacji}
Systemy wykorzystujące LEM POWINNY ograniczać zakres udostępnianych informacji lokalizacyjnych do poziomu wymaganego przez określony cel zastosowania.

Udostępniana reprezentacja informacji lokalizacyjnej MOŻE wykorzystywać ograniczony zakres komponentów modelu, jeżeli zachowana zostaje zgodność z wymaganiami danego profilu oraz interpretacją semantyczną LEM.

Model LEM nie definiuje polityk ujawniania informacji ani mechanizmów kontroli dostępu. Zasady te pozostają odpowiedzialnością implementacji.

\section{Spójność semantyczna rozszerzeń a bezpieczeństwo}
Rozszerzenia wykorzystujące model LEM MUSZĄ zachowywać zgodność z podstawową semantyką modelu (zob. 3.1, 4.5, 7.4).

Rozszerzenie NIE MOŻE wprowadzać nowych znaczeń istniejących komponentów ani powodować niejednoznaczności interpretacji informacji lokalizacyjnej.

Z punktu widzenia bezpieczeństwa, rozszerzenie niespełniające tej zasady stanowi dodatkowe ryzyko: pole rozszerzenia błędnie zinterpretowane jako standardowy komponent LEM może prowadzić implementację do niepoprawnego wniosku o lokalizacji obiektu. Implementacje POWINNY jednoznacznie odróżniać dane rozszerzeń od komponentów podstawowych modelu LEM podczas przetwarzania (zob. 5.1, 7.7).

Dodatkowe informacje wymagane przez określone środowisko zastosowania POWINNY być definiowane poprzez mechanizmy rozszerzeń lub profile, a nie poprzez zmianę podstawowego modelu semantycznego LEM.

\section{Zbiorczy model zagrożeń (informacyjny)}
Poniższa tabela zestawia w jednym miejscu zagrożenia, których dotyczą zasady rozproszone w sekcjach 6.1–6.8, wraz z mechanizmami obrony — pozostającymi, zgodnie z 6.1, poza zakresem modelu LEM. Ma charakter informacyjny i poglądowy, a nie wyczerpującej, formalnej analizy zagrożeń.

\begin{table}[p]
\centering
\footnotesize
\setlength{\tabcolsep}{3pt}
\begin{tabular}{|p{0.35cm}|p{1.9cm}|p{2.3cm}|p{2.5cm}|p{2.5cm}|p{0.75cm}|}
\hline
\textbf{Nr} & \textbf{Zagrożenie} & \textbf{Wektor} & \textbf{Konsekwencja} & \textbf{Mechanizm obrony (poza zakresem LEM)} & \textbf{Odn.} \\ \hline
1 & Podsłuch transmisji & Przechwycenie danych lokalizacyjnych podczas przesyłania między systemami & Ujawnienie lokalizacji osobom nieuprawnionym & Szyfrowanie transportu (np. TLS) & 6.1 \\ \hline
2 & Naruszenie integralności w tranzycie & Modyfikacja wartości komponentów bez wykrycia & Błędna interpretacja lokalizacji — szczególnie krytyczne w profilu lokalizacji awaryjnej (4.3.3) & Podpisywanie, sumy kontrolne & 6.2 \\ \hline
3 & Fałszowanie źródła danych & Wstrzyknięcie danych z podszywającego się lub niezaufanego źródła & Zaufanie do nieprawdziwej lokalizacji & Uwierzytelnianie źródła & 6.4 \\ \hline
4 & Deanonimizacja przez korelację & Połączenie informacji lokalizacyjnej z identyfikatorem zewnętrznym lub wzorcem & Ujawnienie tożsamości osoby mimo formalnej separacji lokalizacja $\leftrightarrow$ identyfikacja & Kontrola dostępu, ograniczenie łączenia zbiorów danych & 6.3, 6.6 \\ \hline
5 & Nadmierne ujawnienie precyzji & Udostępnienie pełnego, precyzyjnego opisu lokalizacji tam, gdzie cel wymagał mniejszej precyzji & Ujawnienie więcej informacji niż niezbędne do celu & Minimalizacja zakresu ujawnianych komponentów & 6.7, 2.4.4 \\ \hline
6 & Wstrzyknięcie nierozpoznanego rozszerzenia & Dodanie pola o nazwie/strukturze imitującej komponent LEM & Błędna interpretacja danych rozszerzenia jako części modelu semantycznego & Jawne, techniczne rozróżnianie rozszerzeń od komponentów podstawowych & 6.8, 5.1 \\ \hline
7 & Obejście wymagań profilu & Świadome przesłanie opisu zgodnego z modelem podstawowym, lecz nie spełniającego wymagań profilu & Działanie na podstawie informacji o niewystarczającej precyzji & Wymuszenie walidacji zgodności z profilem po stronie odbiorcy & 4.4.3, 4.3.3.1 \\ \hline
8 & Wykorzystanie nieaktualnych danych & Odebranie i wykorzystanie przestarzałego stanu informacji lokalizacyjnej & Działanie na podstawie nieaktualnej lokalizacji obiektu & Znaczniki czasu, mechanizmy świeżości danych & A.12, 6.4 \\ \hline
\end{tabular}
\caption{Zbiorczy model zagrożeń (informacyjny)}
\label{tab:threat_model}
\end{table}

Pozycje 6 i 7 wynikają bezpośrednio z właściwości modelu LEM opisanych w innych rozdziałach (odpowiednio: rozróżnianie rozszerzeń od komponentów podstawowych oraz niezależność walidacji profilu od walidacji modelu podstawowego) i ilustrują, że dyscyplina semantyczna modelu ma również praktyczne konsekwencje dla bezpieczeństwa, a nie wyłącznie dla spójności interpretacji.
